\usetikzlibrary{matrix}
\newcommand{\KernelRangeFigure}{%
\begin{tikzpicture}[>=Stealth,font=\small,x=1cm,y=1cm]
\foreach \dx/\kind in {0/1,6/2}{
\begin{scope}[xshift=\dx cm]
 \draw[->,gray](-.5,0)--(2.8,0)node[right,text=black]{$x$};
 \draw[->,gray](0,-.4)--(0,2.65)node[above,text=black]{$y$};
 \draw[blue](-.45,0)--(2.65,0);
 \draw[->,gray](0,0)--(1,2);
 \fill (1,2)circle(1.5pt)node[above right]{$v=(1,2)$};
 \ifnum\kind=1
  \node at (1.15,3.3) {$P(x,y)=(0,y)$};
  \draw[red] (0,-.35)--(0,2.5);
  \draw[red,thick,->](0,0)--(0,2)node[left]{$Pv$};
  \draw[dashed,gray](1,2)--(0,2);
  \node[below,text=blue]at(1.5,-.1){$\ker P$};
  \node[left,text=red]at(0,1){$\im P$};
 \else
  \node at (1.15,3.3) {$N(x,y)=(y,0)$};
  \draw[red,dashed](-.45,0)--(2.65,0);
  \draw[red,thick,->](0,0)--(2,0)node[below]{$Nv$};
  \draw[dashed,gray,->](1,2)--(2,0);
  \node[below]at(1.15,-.65){$\ker N=\im N=\mathbb Re_1$};
 \fi
\end{scope}}
\end{tikzpicture}}

\newcommand{\EvaluationColumnsFigure}{%
\begin{tikzpicture}[font=\small]
\matrix (a) [matrix of math nodes,ampersand replacement=\&,left delimiter=(,right delimiter=),
 nodes={minimum width=9mm,minimum height=5mm,inner sep=1pt},
 column 1/.style={nodes={text=blue}},column 2/.style={nodes={text=red}},
 column 3/.style={nodes={text=green!50!black}}]
 {1\&1\&1\&1\\1\&i\&-1\&-i\\1\&-1\&1\&-1\\1\&-i\&-1\&i\\};
\node[left=5mm of a] {$A=M(T)=$};
\foreach \j/\label in {1/{T1},2/{Tx},3/{Tx^2},4/{Tx^3}}
 \node[above=3mm of a-1-\j] {$\label$};
\end{tikzpicture}}

\newcommand{\MarkovTransitionFigure}{%
\begin{tikzpicture}[>=Stealth,font=\small,
 state/.style={draw,circle,minimum size=7mm}]
\node[state](a) at (0,0){$A$};
\node[state](b) at (3,0){$B$};
\node[state](c) at (6,0){$C$};
\path[->] (a)edge[loop above]node{$1/2$}(a)
 (b)edge[loop above]node{$1/2$}(b)
 (c)edge[loop above]node{$1/2$}(c)
 (a)edge[bend left=20]node[above]{$1/2$}(b)
 (b)edge[bend left=20]node[below]{$1/4$}(a)
 (b)edge[bend left=20]node[above]{$1/4$}(c)
 (c)edge[bend left=20]node[below]{$1/2$}(b);
\end{tikzpicture}}
